\chapter{Implementation}
\label{ch:implementation}

\section{System Architecture}
\label{sec:system-architecture}

The implementation separates dataset preparation, page-level
annotation-geometry analysis, model training, prediction export, and benchmark
evaluation into dedicated scripts under \texttt{code/script}. Shared dataset
paths, model names, clustering defaults, and evaluation thresholds are defined
in \texttt{config.py}; run-specific training settings are supplied explicitly
through command-line arguments in the version-controlled SLURM launch scripts.
This organization supports the staged workflow used in the thesis: data
preparation, Stage~2 analysis, Stage~3 baseline evaluation, and Stage~4
discriminative modeling.

% ---------------------------------------------------------------------------
\section{Data Preprocessing}
\label{sec:data-preprocessing}

The data preprocessing pipeline was extended substantially during the thesis
because the raw KVP10k rows were not directly usable for the Mistral baseline.
In particular, the learned baseline required page-level document text, bounding
boxes, and aligned key-value pair supervision in a format consistent with the
intended IBM-style prompt and target representation.

\subsection{Image Processing}
\label{sec:image-processing}

For the Mistral baseline, preprocessing begins from the original PDF source
referenced by each dataset item. The relevant page is accessed using the
dataset page number and converted from the dataset's one-indexed page numbering
to PyMuPDF's zero-indexed page access. Page size is preserved and later mapped
to a quantised $100 \times 100$ coordinate system for prompt and target
generation.

Although the Stage~4 data class supports loading rendered page images, the
reported Stage~4 runs did not provide document-image content. When no image is
loaded, the data class supplies a blank white $224 \times 224$ placeholder to
the LayoutLMv3 processor. Final V4 was launched without
\texttt{--include\_images}; consequently, its document-specific inputs were the
extracted text tokens and their bounding boxes, not page pixels.

\subsection{Text Processing}
\label{sec:text-processing}

Instead of relying on an external OCR engine, the implementation uses PyMuPDF
to extract word-level text and bounding boxes directly from the PDF text layer.
This choice was made for practical deployment reasons on the cluster
environment, where adding a system OCR dependency would have complicated
execution. The trade-off is that scanned pages without a text layer cannot
contribute usable words and are therefore discarded during preparation.

The extracted words are grouped into line-like units and converted into an
LMDX-style textual document representation. This representation is later
embedded into the prompt given to the Mistral model. Ground-truth targets are
serialised as Python list-of-lists strings, following the benchmark style used
by the Stage 3 pipeline.

\paragraph{V4: Word-level splitting.}
Diagnostic analysis (Section~\ref{sec:exp567}) revealed that
the LMDX text format stores entire OCR text blocks --- often multi-word
phrases such as \texttt{1~lb.\ shelled \$15.00} --- as single ``words''
with a single bounding box. This is problematic for LayoutLMv3, which expects
individual words, and for ground-truth label generation, because a key-value
pair whose key and value text reside within the same OCR block cannot produce
distinct entity labels or link labels at the block level.

The V4 data pipeline splits each OCR block into whitespace-separated words and
estimates horizontal sub-bounding boxes proportional to character count. For
example, the block \texttt{1~lb.\ shelled \$15.00} with bounding box
$[x_1, y_1, x_2, y_2]$ is split into four words (\texttt{1},
\texttt{lb.}, \texttt{shelled}, \texttt{\$15.00}), each assigned a
horizontal slice of the original bbox. This increased the average word count
per document from $\sim$35 (block-level) to $\sim$200 (word-level) and
enabled the entity labels and link labels to reach word-level granularity.

\subsection{Layout Encoding}
\label{sec:layout-encoding}

Annotation geometry is fused with extracted words through an OCR-to-annotation
matching step. Bounding boxes are converted into normalised integer coordinates
on a $100 \times 100$ grid. The resulting aligned entities are then transformed
into three categories: regular key-value pairs, unkeyed values, and unvalued
keys. This preprocessing stage also handles repeated annotator copies in
KVP10k by grouping rows by page identifier and selecting the richest annotation
set for each unique page.

For Stage~4, bounding box coordinates are further re-normalised to the
$1000 \times 1000$ LayoutLMv3 convention (integers in $[0, 1000]$) after the
above deduplication step.

% ---------------------------------------------------------------------------
\section{Model Implementation}
\label{sec:model-implementation}

\subsection{LayoutLMv3 Integration}
\label{sec:layoutlmv3-integration}

Stage~4 uses \texttt{LayoutLMv3KVPModel} for V1 and
\texttt{LayoutLMv3KVPModelV2} for V2/V4; both wrap the
\texttt{microsoft/layoutlmv3-base} checkpoint from the HuggingFace Hub.
The base encoder is loaded in full fp32 precision. No quantisation is applied:
the 125M-parameter LayoutLMv3-base model fits within the 40\,GB memory
envelope of the A100 nodes used for Stage~4, unlike the 7B-parameter Mistral
model that required QLoRA in Stage~3. LayoutLMv3 remains a multimodal
architecture, but the blank visual placeholders used here carry no
document-specific image information.

All encoder parameters are unfrozen and participate in fine-tuning. A linear
learning-rate warmup of 500 steps is used to avoid large gradient updates on
the pre-trained weights during the first portion of training (see
Section~\ref{sec:optimisation}).

Tokenisation follows the LayoutLMv3 tokenizer convention: each OCR word and
its bounding box are passed together, and sub-word tokens inherit the bounding
box of their parent word. The maximum sequence length is capped at 512 tokens;
longer pages are truncated from the right after the \texttt{[SEP]} position.

\subsection{Entity Classifier Head}
\label{sec:entity-classifier-impl}

The entity classification head is a two-layer MLP applied independently to
each text token output of the encoder (image patch tokens are excluded from
classification). The network maps the 768-dimensional token representation to
a 3-class logit vector (\texttt{O}, \texttt{KEY}, \texttt{VALUE}) as described
in Section~\ref{sec:entity-classifier}.

During training, the cross-entropy loss is computed over all text tokens.
Tokens corresponding to padding positions and to image patches are masked
from the loss. To address the class imbalance (the \texttt{O} class typically
represents $\sim$70\% of tokens), class weights are set to the inverse
frequency observed across a 1{,}000-sample calibration pass of the training
set at the start of each Stage~4a training run.

\subsection{Biaffine Linker Module}
\label{sec:linker-impl}

The linker module receives key and value representations extracted from the
entity classifier output. In V1, these are individual token hidden states;
in V2, they are span-level mean-pooled representations
(Section~\ref{sec:span-representation}).

Representations are augmented with 8 spatial features
(Section~\ref{sec:spatial-encoding}) before scoring. The semantic scorer uses
scaled dot-product attention (Section~\ref{sec:kv-attention}) --- an earlier
bilinear implementation was abandoned due to gradient explosion.

\paragraph{Teacher forcing.}
During training, the linker receives ground-truth entity labels for span
grouping rather than the entity classifier's \texttt{argmax} predictions.
This ensures the linker learns from correctly segmented spans. At inference
time, the linker receives predicted entity labels. This teacher-forcing
strategy decouples entity classification errors from linker training and is
standard practice in structured prediction~\cite{spade}.

\paragraph{NaN guard.}
Achieving stable training required a valid-label mask before the BCE call:
ground-truth link labels outside $[0, 1]$ (from uninitialised positions in
the adjacency matrix) are excluded.  This is analogous to the
\texttt{ignore\_index=-100} mechanism used for entity labels in
\texttt{CrossEntropyLoss}. The \texttt{pos\_weight} for the BCE loss is set
to $\min(n_{\text{neg}}/n_{\text{pos}},\, 50)$ in V1 and
$n_{\text{neg}}/n_{\text{pos}}$ (uncapped) in V2.

\paragraph{Memory efficiency.}
To avoid materialising large expanded tensors, scoring is performed in chunks
of 8 key rows. A hard cap of 128 key and 128 value candidates per page
(selected by highest entity-classifier confidence) bounds the worst-case
linker matrix to $128 \times 128$ in V1.

% ---------------------------------------------------------------------------
\section{Baseline Evaluation Pipeline}
\label{sec:baseline-eval-pipeline}

Stage~3 converts KVP10k annotations and generated predictions into the released
benchmark format, then applies type-aware greedy matching in text-only,
location-only, and combined modes. Precision and recall are averaged per page
before F1 is computed. The evaluation used 581 prepared unique test pages; the
5{,}255 rows in the raw test split are duplicated annotation copies and are not
the evaluation sample count.

% ---------------------------------------------------------------------------
\section{Training Pipeline}
\label{sec:training-pipeline}

The training pipeline for the Mistral baseline (Stage~3) operates on the
prepared JSON files generated during preprocessing. Each sample contains the
prompt text, the expected target response, and the aligned ground-truth
key-value pairs. During tokenisation, prompt tokens are masked in the loss so
that optimisation is applied only to the response segment.

\subsection{Data Loader}
\label{sec:data-loader}

The data loader reads page-level JSON files from the prepared train and test
directories. Samples without prompt text or without non-empty supervision are
skipped to avoid training on degenerate examples. For Stage~4, a
\texttt{LayoutLMv3PreparedDataset} wraps the same JSON files, parses words and
bounding boxes from the LMDX representation, creates entity and link labels,
and tokenizes them to a maximum length of 512. The reported runs used the
loader's blank visual placeholder rather than a page-image cache.

\subsection{Training Loop}
\label{sec:stage4-training-loop}

\paragraph{Stage~3 (Mistral QLoRA).}
We used QLoRA with 4-bit NF4 quantisation due to memory constraints on the
available 40\,GB A100 GPU.

\paragraph{Stage~4a -- Entity Classifier Pre-training.}
Stage~4a trains the LayoutLMv3 encoder together with the entity classification
head only, with the linker module disabled ($\lambda = 0$).
This pre-trains the representation before the more complex
linking objective is introduced. Training uses fp32 precision
on an A100 node; no quantisation is required. A per-device batch size of 4 with
8 gradient-accumulation steps gives an effective batch size of 32. The Stage~4a
run uses early-stopping patience 3 based on validation entity F1.

\paragraph{Stage~4b -- Joint Fine-tuning.}
Stage~4b activates the biaffine linker ($\lambda > 0$). Earlier V1 and V2
development runs were initialized from the Stage~4a entity checkpoint. Final
V4, which uses the V2 span-linking architecture with the data-pipeline fixes
below, was instead warm-started from
\texttt{stage4b\_canary\_B/best\_model/pytorch\_model.bin}. It used a
per-device batch size of 1, 8 gradient-accumulation steps (effective batch size
8), a unified learning rate of $2\times10^{-5}$, linker-loss weight
$\lambda=5$, and early-stopping patience 10 based on validation entity F1.
Training targeted 30 epochs across resumed jobs rather than one uninterrupted
30-epoch run.

\paragraph{V4 link label generation.}
Diagnostic analysis of V2 training revealed that the ground-truth link label
matrix contained a cross-product explosion: for each key-value pair, \emph{all}
key-word indices were linked to \emph{all} value-word indices, rather than a
single representative pair. A document with 8 ground-truth KVPs produced
$\sim$1{,}049 positive link entries instead of~8, causing the span-level
collapse (Section~\ref{sec:v2-span-linking}) to mark nearly every span pair as
positive --- equivalent to no supervision at all.

The V4 fix assigns exactly one link per KVP: the first key word is linked to
the first value word, producing $N$ positive link entries for $N$ ground-truth
KVPs.  Together with the word-level splitting
(Section~\ref{sec:text-processing}), a bidirectional bounding-box overlap
criterion ($\geq 50\%$ of either the word or the target region), and a relaxed
span grouping filter that only excludes $[0,0,0,0]$ bounding boxes (previously,
the filter $x_1 < x_2 \wedge y_1 < y_2$ rejected $\sim$50\% of valid tokens
whose LayoutLMv3-normalised coordinates had $x_1 = x_2$), the V4 pipeline
recovers 141 of 179 word-level links (79\%) and 125 of 179 span-level links
(70\%) in the first 20 prepared test pages used for the diagnostic in
Section~\ref{sec:res567}. These are label-pipeline recovery checks, not
full-test model-performance metrics.

\subsection{Checkpointing}
\label{sec:checkpointing}

Model checkpoints are stored after each epoch and the best checkpoint (highest
validation entity F1) is retained. For Stage~4b, the serialized
\texttt{pytorch\_model.bin} state dictionary contains the encoder, entity
classifier, and linker parameters. Prediction export and benchmark evaluation
then load this fixed checkpoint independently of training.

% ---------------------------------------------------------------------------
\section{Evaluation Protocol}
\label{sec:evaluation-protocol}

We implement a dual-threshold evaluation protocol:

\subsection{Text Matching: NED}
\label{sec:ned}

Normalised Edit Distance (NED) measures text accuracy:
\[
  \text{NED}(p, t) = \frac{\text{Levenshtein}(p, t)}{\max(|p|, |t|)}
  \tag{5.1}
\]
where $p$ is predicted text, $t$ is target text. Threshold: $\text{NED} < 0.2$.

\subsection{Location Matching: IoU}
\label{sec:iou}

Intersection over Union (IoU) measures spatial accuracy:
\[
  \text{IoU}(B_p, B_t) = \frac{\text{Area}(B_p \cap B_t)}{\text{Area}(B_p \cup B_t)}
  \tag{5.2}
\]
where $B_p$ is predicted box and $B_t$ is target box. The paper states
$\text{IoU}>0.3$, but IBM's released evaluator checks
$\text{IoU}\geq0.3$; reported benchmark results follow the released code.

\subsection{Overall Metric: F1 Score}
\label{sec:f1}

A regular pair matches only if both key and value satisfy
$\text{NED}<0.2$ and, for combined evaluation, $\text{IoU}\geq0.3$.
Predicted and ground-truth KVP types must agree, and greedy matching assigns
each ground-truth item at most once. For each scored page $i$, precision $P_i$
and recall $R_i$ are computed from pair-level counts. The released evaluator
then reports
\[
  \bar P = \frac{1}{D}\sum_{i=1}^{D} P_i, \qquad
  \bar R = \frac{1}{D}\sum_{i=1}^{D} R_i, \qquad
  F_1 = \frac{2\bar P\bar R}{\bar P+\bar R},
  \tag{5.3}
\]
where $D$ contains pages with non-empty ground truth after category filtering.
Text-only, location-only, and combined metrics are reported separately for
regular, unkeyed, unvalued, and All categories.

% ---------------------------------------------------------------------------
\section{Code Organisation}
\label{sec:code-organisation}

The principal implementation files are:
\begin{verbatim}
code/script/
|-- prepare_data.py              # PDF text and KVP preparation
|-- stage4_kvp_dataset.py        # Stage 4 dataset and labels
|-- layoutlm_model.py            # V1 token linker
|-- layoutlm_model_v2.py         # V2/V4 span linker
|-- train_stage4a.py             # Entity-only training
|-- train_stage4b_v2.py          # V2/V4 joint training
|-- export_v4_predictions.py     # Fixed-checkpoint inference
'-- evaluate_kvp10k_benchmark.py # Released benchmark semantics
\end{verbatim}

% ---------------------------------------------------------------------------
\section{Reproducibility}
\label{sec:reproducibility}

Reproducibility is supported by seeded data splitting, explicit CLI arguments
in version-controlled SLURM scripts, per-epoch checkpoints, run logs, and
offline HuggingFace caches on compute nodes. Data preparation and GPU training
remain separate steps because source-PDF retrieval is performed before jobs are
submitted to the offline compute environment.
